% Appendix A: Alloy Formal Models

\section{Alloy Formal Models}
\label{app:alloy}

The bounded model checking referenced in Section~\ref{sec:discussion}
uses Alloy 6 with a three-layer module hierarchy: a Bitcoin base
layer (\texttt{btc\_base}), a vault abstraction layer
(\texttt{vault\_base}), and five concrete covenant models extending
the abstraction with design-specific transition guards. Two
cross-cutting modules (\texttt{threat\_model},
\texttt{cross\_covenant}) are opened by the concrete models. In
total the corpus contains 42 assertions; all but one pass under the
configured scopes.

\subsection{Module Contents}

\paragraph{\texttt{btc\_base.als}.}
Signatures for \texttt{UTXO}, \texttt{Key}, \texttt{Transaction},
and \texttt{Address}. UTXOs carry an amount, owner key, and
spent-status flag; transactions consume inputs and produce outputs.

\paragraph{\texttt{vault\_base.als}.}
Abstract vault state machine with four states (\texttt{VAULTED},
\texttt{UNVAULTING}, \texttt{WITHDRAWN}, \texttt{RECOVERED}),
transitions (trigger, withdraw, recover), and a \texttt{VaultFamily}
relation grouping UTXOs of a single vault instance. The CSV delay
is an integer field on each vault UTXO. Invariants defined at this
level (fund conservation, recovery liveness) propagate to all
concrete models.

\paragraph{Concrete vault models.}
Five modules extend \texttt{vault\_base} with covenant-specific
transition guards: \texttt{ctv\_vault.als} models the template-hash
commitment; \texttt{ccv\_vault.als} models the mode enumeration and
\texttt{OP\_SUCCESS} behaviour for undefined modes;
\texttt{opvault\_vault.als} models the three-key separation and
authorised recovery; \texttt{cat\_csfs\_vault.als} models the
dual-verification binding and the unconstrained recovery leaf;
\texttt{simplicity\_vault.als} models the
\texttt{jet::outputs\_hash()} check on both trigger and recovery
paths.

\paragraph{\texttt{threat\_model.als}.}
Attacker profiles parameterised by which keys the attacker
controls; each property is checked under each profile to identify
the minimum capability needed for each attack class.

\paragraph{\texttt{cross\_covenant.als}.}
Models concurrent vaults from different covenant types sharing the
same UTXO set; checks cross-vault injection and fee-wallet
contention between concurrent \opvault{} operations.

\subsection{Asserted Invariants}
\label{app:alloy:invariants}

The main shared invariants are:

\begin{description}[leftmargin=1.4em,itemsep=1pt,parsep=0pt,topsep=2pt]
  \item[Fund conservation.] Total vault value is preserved across
    lifecycle transitions except via recovery or withdrawal.
  \item[Destination binding.] The withdraw path produces only the
    pre-committed withdrawal output (\opctv{}, \opvault{}, Simplicity);
    for \opccv{} the withdraw path is \opctv{}-committed; \catcsfs{}'s hot
    leaf is bound via dual verification.
  \item[Recovery authorisation.] Recovery can only be invoked by a
    holder of the design's authorisation set
    ($\emptyset$ for \opccv{}; cold key for \opctv{}, \catcsfs{}, Simplicity;
    \texttt{recoveryauth} key for \opvault{}).
  \item[Recovery liveness.] From \texttt{VAULTED} or
    \texttt{UNVAULTING}, some transition to \texttt{RECOVERED}
    exists under the authorisation set.
  \item[No cross-leaf extraction.] Spending the trigger leaf
    cannot produce outputs that would satisfy the recovery leaf
    and vice versa.
\end{description}

\subsection{Per-Model Verification Results}

\paragraph{\opctv{}.} Six assertions all pass, including
no-extraction under each single-key compromise profile, CSV
enforcement, and the no-state-proliferation property. The
address-reuse scenario finds an instance where a second deposit
produces an unspendable UTXO, confirming TM7 as structural.

\paragraph{\opccv{}.} Eight assertions pass, including a
mode-confusion containment check (undefined modes trigger
\texttt{OP\_SUCCESS} but only within the affected leaf; the
vault's other leaves remain secure). The griefing scenario
demonstrates that any party may invoke recovery without keys
(TM2). The splitting scenario exhibits the revault chain that
produces the per-UTXO recovery-cost tail of class \classscale{}
(TM4).

\paragraph{\opvault{}.} Eight assertions pass, including a
dual-key-no-theft property: joint compromise of the trigger and
\texttt{recoveryauth} keys yields liveness denial but not theft,
because the recovery address is pre-committed in the Taproot
tweak. A key-loss scenario exhibits permanent inability to recover
if the \texttt{recoveryauth} key is destroyed.

\paragraph{\catcsfs{}.} Five of six assertions pass. The
\texttt{catcsfsNoExtraction\_ColdKeyOnly} assertion \emph{fails}:
Alloy exhibits a trace in which the cold-key holder invokes
recovery with an attacker-chosen output address, extracting funds.
This is the structural confirmation of TM11 and the only expected
assertion failure in the corpus. Destination-lock and CSV
assertions pass.

\paragraph{Simplicity.} Six assertions pass: no-extraction under
any single-key compromise, output binding enforcement on both
trigger and recovery, CSV enforcement, and no state proliferation.
A recovery-constrained scenario confirms that cold-key compromise
can only recover to the pre-committed address, unlike \catcsfs{}.

\subsection{Scope and Limitations}

The Alloy models operate at the state-machine level, abstracting
cryptographic primitives (signatures are modeled as key-possession
checks; hash functions are assumed collision-resistant),
serialisation, and temporal dynamics. Network propagation, mining,
and mempool policy are out of scope. Verification confirms
\emph{structural} reachability and authorisation properties but
not cryptographic unforgeability or temporal race outcomes; the
regtest experiments (Section~\ref{sec:empirical}) complement the
formal layer with measured transaction sizes and empirical attack
costs.

\providecommand{\todo}[2][]{\textbf{[TODO: #2]}}
